% titel: Blatt 9, A9: NO-Bildung im Gleichgewicht
% Teil II: der durchgerechnete Weg.

\subsection*{Gegeben und gesucht}
$N_2 + O_2 \rightleftharpoons 2NO$, $K_c = 0.55$, Start je $2\ mol\,L^{-1}$.
Gesucht: $[NO]$ im Gleichgewicht.

\subsection*{Ansatz}
$[N_2]=[O_2] = 2-x$, $[NO] = 2x$ (in mol/L).

\subsection*{Rechnung}
\[
  K_c = \frac{(2x)^2}{(2-x)^2} = 0.55 .
\]
Wurzel:
\[
  \frac{2x}{2-x} = \sqrt{0.55} \approx 0.7416 .
\]
\[
  2x = 0.7416(2-x)
  \;\Rightarrow\; 2x = 1.4832 - 0.7416 x
  \;\Rightarrow\; 2.7416 x = 1.4832 .
\]
\[
  x \approx 0.541,\qquad [NO] = 2x \approx 1.08\ mol\,L^{-1}.
\]

\begin{ergebnis}
$[NO] \approx 1.08\ mol\,L^{-1}$; dann bleiben $[N_2]=[O_2]\approx1.46$
mol$^{-1}$.
\end{ergebnis}

\begin{kontrolle}
Probe: $\frac{(2\cdot0.541)^2}{(2-0.541)^2}
= \frac{1.171}{2.129} \approx 0.55 = K_c$. Passt. Da $K_c<1$ ist das
Gleichgewicht eher auf der Eduktseite, aber wegen der kleinen
Stöchiometrie bleibt noch genug NO. (Zwischenkonzentrationen sind bulundet,
$2-x>0$ ist erfüllt.)
\end{kontrolle}

\begin{merksatz}
MWG nach $x$ auflösen; $[NO]$ aus $2x$. Wurzelziehen vereinfacht, wenn
Zähler/Nenner gleiche Exponenten haben.
\end{merksatz}
